% GERADO por scripts/slr_pipeline.py a partir de docs/slr/screening.csv
% NÃO editar à mão: os números vêm da triagem real. Ver docs/PROTOCOLO_SLR.md
\begin{figure}[H]
    \centering
    \begin{tikzpicture}[node distance=1.2cm]
    % Sem hifenização dentro das caixas: num texto de 5-6 cm ela produzia
    % «Registos identifica-dos», «tria-dos» e «ve-nue», que num diagrama se leem
    % como gralhas. As caixas são ragged, por isso proibir a hifenização só faz
    % as linhas quebrarem mais cedo — nenhuma palavra excede a largura.
    \hyphenpenalty=10000 \exhyphenpenalty=10000
    \tikzstyle{process} = [rectangle, minimum width=6cm, minimum height=1.1cm,
        text centered, draw=black, fill=white, text width=6cm, font=\small]
    \tikzstyle{saida} = [rectangle, minimum height=1.1cm, draw=black, fill=black!3,
        text width=5.8cm, font=\scriptsize, align=left]
    \tikzstyle{arrow} = [thick,->,>=stealth]
    \tikzstyle{lado} = [rotate=90, anchor=south, font=\bfseries\footnotesize,
        color=black!60]

    \node (id) [process] {Registos identificados nas bases de dados \\
        \scriptsize (IEEE 427 + Scopus 456) \\ \textbf{(n = 883)}};
    \node (dup) [saida, right=1.3cm of id] {Duplicados removidos \\
        (DOI ou título normalizado) \\ \textbf{(n = 203)}};

    \node (screen) [process, below=1.2cm of id] {Registos únicos triados
        (título e resumo) \\ \textbf{(n = 680)}};
    \node (exc1) [saida, right=1.3cm of screen] {Excluídos na triagem
        \textbf{(n = 201)} \\ $\bullet$ Robô único (sem enxame): 75 \\ $\bullet$ Revisão / sem validação experimental: 42 \\ $\bullet$ Sem dados quantitativos: 1 \\ $\bullet$ Sem revisão por pares / integridade: 4 \\ $\bullet$ Duplicado não detetado automaticamente: 4 \\ $\bullet$ Método fora do âmbito (nem MARL nem bio-inspirado): 75};

    \node (full) [process, below=3.5cm of screen] {Avaliados em texto integral \\
        \textbf{(n = 479)}};
    \node (exc2) [saida, right=1.3cm of full] {Excluídos na elegibilidade
        \textbf{(n = 421)} \\ $\bullet$ Tarefa fora do âmbito (alocação, SLAM, comunicações, manipulação): 79 \\ $\bullet$ Fora do núcleo (não é enxame em venue de referência): 342};

    \node (inc) [process, below=2.5cm of full, fill=black!5]
        {\textbf{Estudos incluídos na revisão} \\ \textbf{(n = 58)}};

    \draw [arrow] (id) -- (screen);
    \draw [arrow] (id) -- (dup);
    \draw [arrow] (screen) -- (full);
    \draw [arrow] (screen) -- (exc1);
    \draw [arrow] (full) -- (inc);
    \draw [arrow] (full) -- (exc2);

    % As quatro fases ANCORADAS às caixas, e não a alturas fixas. Com as
    % coordenadas escritas à mão, a «Elegibilidade» e a «Inclusão» ficavam cada
    % uma um degrau acima da caixa a que pertencem e a caixa final acabava sem
    % etiqueta nenhuma — o `(col |- nó)` dá o x da coluna e o y do nó, pelo que
    % a etiqueta segue a caixa mesmo quando o número de motivos de exclusão
    % muda a altura do diagrama.
    \coordinate (col) at (-4.2, 0);
    \node [lado] at (col |- id) {Identificação};
    \node [lado] at (col |- screen) {Triagem};
    \node [lado] at (col |- full) {Elegibilidade};
    \node [lado] at (col |- inc) {Inclusão};
    \end{tikzpicture}
    \caption[Fluxograma PRISMA 2020 da revisão conduzida]{Fluxograma PRISMA 2020 da
    revisão conduzida. Os números são gerados automaticamente a partir do registo de
    triagem (\texttt{docs/slr/screening.csv} no repositório, Anexo~\ref{apx:resources}),
    onde cada decisão está associada ao respetivo critério de exclusão.}
    \label{fig:prisma}
\end{figure}
